\section{Introduction}

Hierarchical structure is one of the recurring reasons that Euclidean representation learning becomes inefficient. Taxonomies branch exponentially with depth, lexical relations mix tree-like and symmetric components, knowledge graphs combine multi-relational hierarchies with many-to-many relations, and graph or language representations often need to place coarse categories, fine instances, and lateral neighborhoods in the same embedding system. Hyperbolic representation learning addresses this mismatch by using negative curvature, whose volume growth matches trees more naturally than Euclidean space \citep{nickel2017poincare,nickel2018lorentz,chami2019hgcn,tifrea2019poincare,balazevic2019poincare}. In examples such as WordNet taxonomies, Freebase relation graphs, and synthetic trees with heterogeneous branching, the relevant question is not only whether hyperbolic space helps, but which curvature each structural component should use during optimization.

\begin{table*}[t]
\centering
\caption{\textbf{HyperbolicMomentum update.} Each group has a dedicated Lorentz manifold. The main path updates the point and tangent momentum; the slower branch changes curvature and maps both states. The curvature derivative is total at fixed canonical point, as in Eq.~\eqref{eq:total-curv-gradient}.}
\label{tab:method-overview}
\footnotesize\setlength{\tabcolsep}{4.2pt}\renewcommand{\arraystretch}{1.12}
\begin{tabularx}{\linewidth}{l l X}
\toprule
Stage & State or operation & Geometric role \\
\midrule
\multicolumn{3}{l}{\textit{Tangent-state and geodesic update}} \\
Gradient & $\widehat g_{g,t}=\grad_{x_g}\mathcal J$ & Form a tangent direction at the current point. \\
Momentum & $\widetilde m_{g,t}=\beta m_{g,t}+(1-\beta)\widehat g_{g,t}$ & Average directions in the same tangent space. \\
\rowcolor{oursbg}
Normalized step & $\eta^{\rm geo}_{g,t}=\eta_t/(1+\gamma\sqrt{K_{g,t}}\|\widetilde m_{g,t}\|)$ & Control the curvature--momentum scale. \\
Geodesic move & $x_{g,t+1/2}=\Exp_{x_{g,t}}^{K_{g,t}}(-\eta^{\rm geo}_{g,t}\widetilde m_{g,t})$ & Move the point on its current manifold. \\
Tangent transport & $m_{g,t+1/2}=\PT_{x_{g,t}\to x_{g,t+1/2}}^{K_{g,t}}\widetilde m_{g,t}$ & Carry memory to the new base point. \\
\midrule
\multicolumn{3}{l}{\textit{Slow curvature update and compatible joint-state map}} \\
\rowcolor{oursbg}
Learn curvature & $\widehat h_{g,t}=\partial_{c_g}F(q_t,u_t)$; Eq.~\eqref{eq:k-step-main} & Use the total derivative, with bounded log drift in implementation. \\
\rowcolor{oursbg}
Map joint state & $a_g=\sqrt{K_{g,t}/K_{g,t+1}}$; Eq.~\eqref{eq:xcurv-main} & Scale both point and momentum by $a_g$; restore tangency numerically. \\
\bottomrule
\end{tabularx}
\end{table*}


Table~\ref{tab:method-overview} summarizes the coupled update. A useful distinction is between fixed-manifold optimization, architectural curvature learning, and consistency of optimizer state when curvature changes. Recent curvature-aware Lorentz optimization already addresses the last issue using origin-tangent remapping~\citep{bdeir2025robust}. Our study concerns the particular homothetic state map, grouped slow curvature updates, normalized geodesic step, and their conditional analysis.


A fixed-curvature optimizer cannot adapt its manifold to different structural groups. Riemannian stochastic gradient descent provides a principled first-order update on manifolds \citep{bonnabel2013sgd,zhang2016geodesic}, and Riemannian adaptive methods extend adaptive optimization to curved spaces \citep{becigneul2019radam}. Yet a fixed-curvature optimizer cannot adapt when one relation behaves like a taxonomic branch, another relation behaves like a symmetric neighborhood, and a third relation behaves like a high-degree multi-relational association. Knowledge graph baselines such as TransE and RotatE illustrate how strongly relation type shapes the geometry of a representation problem \citep{bordes2013transe,sun2019rotate}; hyperbolic KG methods further show that relation-aware curvature can matter \citep{balazevic2019poincare}. A second requirement is geometric: momentum vectors live in tangent spaces. The vector carried from step $t-1$ belongs to $T_{x_{t-1}}\HH^d_{-K}$, whereas the current gradient belongs to $T_{x_t}\HH^d_{-K}$, so adding them without parallel transport has no intrinsic meaning. If $K$ itself changes, then both the base point and the metric radius change, making a point-only curvature update incomplete for architectures in which each group has its own manifold.

This paper studies the optimization problem that appears when these two requirements coincide. We ask whether a hierarchical model should learn \textbf{one curvature per structural group}, whether the geodesic step should be normalized by the local curvature-momentum scale, and when a \textbf{cross-curvature momentum state} is actually meaningful. The answer is deliberately nuanced. Historical unfiltered, tail-only experiments with Euclidean Adam provide exploratory evidence for grouped learnable curvature: on WN18RR, B4 improves Hits@10 by $+4.1\%$ and loss by $-2.8\%$ over fixed curvature B3, while MRR is positive but sensitive to the moving-average window. On FB15k-237, the same mechanism gives a different metric profile: MRR improves by $+2.2\%$ and Hits@1 by $+3.1\%$, but Hits@10 is almost unchanged at $+0.1\%$. Thus the shared positive KG signal is MRR, while which Hits metric benefits depends on the dataset.

We present \textbf{HyperbolicMomentum}, an optimizer in the Lorentz model with three mechanisms: grouped learnable curvature, curvature-normalized geodesic steps, and cross-curvature parallel-transport momentum. The method is motivated by one controlling geometric quantity. The product $\sqrt{K}R$ governs capacity, conditioning, and stability over radius $R$, while $\sqrt{K}\eta\|m\|$ governs exponential-map distortion during an optimization step. Our theory makes this trade-off explicit through exact homothety and Jacobi-field identities, a quantitative heterogeneity bound, error-explicit two-block convergence, momentum tracking, constrained curvature stationarity, finite-precision propagation, and augmented-state stability with an explicit training-time factor. Table~\ref{tab:curvatures} lists the relation-dependent curvature states: WN18RR learns a wide relation-dependent curvature spread, whereas FB15k-237 learns a much smaller spread consistent with weaker hierarchy.

\begin{table*}[t]
\centering
\caption{\textbf{Learned relation curvatures.} Historical B4 values. All 11 WN18RR relations and the 15 displayed FB15k-237 relations are retained; the fixed baseline is $K=1$. Curvature magnitude is a parameter, not a performance ranking.}
\label{tab:curvatures}
\footnotesize\setlength{\tabcolsep}{4.2pt}\renewcommand{\arraystretch}{1.12}
\begin{adjustbox}{max width=\linewidth}
\begin{tabular}{lrrr}
\toprule
WN18RR relation & $K$ & FB15k-237 relation & $K$ \\
\midrule
derivationally related form & 2.0000 & award/nominee & 1.8371 \\
member meronym & 1.5823 & film/release region & 1.3486 \\
hypernym & 1.4173 & award/winner (A) & 1.3449 \\
has part & 1.3050 & award/nominee$\to$award & 1.2725 \\
instance hypernym & 1.2220 & music/genre$\to$artists & 1.1832 \\
verb group & 1.0938 & person/profession & 1.1817 \\
also see & 1.0888 & film/genre & 1.1754 \\
synset domain topic of & 1.0754 & award/nominees for & 1.1678 \\
member of domain region & 1.0170 & actor/filmography & 1.1677 \\
member of domain usage & 1.0092 & award/winners (B) & 1.1435 \\
similar to & 1.0084 & film/other crew & 1.1421 \\
-- & -- & award/nominated for & 1.1338 \\
-- & -- & person/nationality & 1.1093 \\
-- & -- & work/awards won & 1.1034 \\
-- & -- & sports/team roster & 1.0966 \\
\bottomrule
\end{tabular}
\end{adjustbox}
\end{table*}


The contributions are as follows, stated in prose to match the empirical evidence. First, we formulate HyperbolicMomentum as a single Lorentz-model algorithm that learns grouped curvatures and uses curvature-normalized geodesic steps, with cross-curvature momentum transport included only when the architecture supplies a dedicated manifold per group. Second, we derive an exact curvature-independent sensitivity certificate for normalized geodesic steps, optimal normalization derivative constants, and a sharp Hessian-weighted heterogeneity bound. The convergence proof preserves finite-horizon martingale weights, while the statistical analysis combines the full state-comparison matrix with stability concentration and entropy chaining. Third, we separate historical curvature ablations from new full-data comparisons and 2000-update controlled drift experiments. The complete matched-initialization search yields held-out distortion reductions of $2.94\%$, $3.31\%$, and $5.02\%$ relative to selected normalized fixed curvature in the stationary, $2\times$ smooth-drift, and $8\times$ smooth-drift conditions. These are the stronger conditions in a six-case matrix; the point-only transport ablation remains close. These conditional gains combine scale adaptation with the physical-distance budget under a canonical-radius cap; they do not imply a uniform advantage across all conditions. Direct operator measurements separately verify the sensitivity and momentum-consistency predictions. Fourth, we document structural curvature differentiation: WN18RR differentiates across relations, FB15k-237 by relation complexity with a four-times smaller $K$ spread, and synthetic trees by branching rate.
